\section{What a fixed-reasoning gradient can and cannot see}
\label{sec:theory}

An edit $P\to P'$ changes accuracy (Eq.~\ref{eq:J}) through two channels: the reasoning the model
writes, and how an answer is read off a given reasoning. With the likelihood ratio of a reasoning
under the two prompts, $w(r)=\pi_{P'}(r\mid x)/\pi_P(r\mid x)$, and $\mathbb{E}_{\pi_P}[w\,g]=\mathbb{E}_{\pi_{P'}}[g]$,
\begin{equation}
J(P')-J(P)=\underbrace{\mathbb{E}_{x}\mathbb{E}_{r\sim\pi_P}\big[(w(r)-1)\,c_{P'}(x,r)\big]}_{D\text{: reasoning distribution}}
+\underbrace{\mathbb{E}_{x}\mathbb{E}_{r\sim\pi_P}\big[c_{P'}(x,r)-c_P(x,r)\big]}_{A\text{: answer read-off}} .
\label{eq:decomp}
\end{equation}
This identity is exact (for $\tau>0$, where every reasoning has positive probability under both
prompts). \greater{}'s objective (Eq.~\ref{eq:greater}) is a soft version of the read-off term $A$
evaluated on a single reasoning, the incumbent's greedy trace, and its gradient is the first-order
change of that quantity; it carries no information about $D$. To first order in a small edit,
$D\approx\operatorname{Cov}_{\pi_P}(\log w,\,c)$: an edit helps if it makes the model's own
\emph{correct} reasoning more likely relative to its incorrect reasoning---the score-function
term that a gradient taken over a fixed reasoning drops.

\begin{proposition}[Faithful reading]
\label{prop:faithful}
If, for every candidate $P'$, the answer read under $P'$ equals the answer read under $P$ except
on a set of reasonings of $\pi_P$-probability at most $\varepsilon$, then $|A|\le\varepsilon$ and
$J(P')-J(P)=D\pm\varepsilon$. Moreover, on a question whose reasoning $r$ supports a wrong answer
$a\neq y$, every decrease of \greater{}'s loss on $r$ raises the probability of reading an answer
that the reasoning does not support, and a decrease by $\delta>0$ nats implies
$q_{P'}(a\mid x,r)\le1-e^{\delta}q_P(y\mid x,r)$.
\end{proposition}

The proof is immediate from Eq.~\ref{eq:decomp} and $q_{P'}(a\mid x,r)+q_{P'}(y\mid x,r)\le1$.
The regime is not hypothetical. Measured directly on the incumbents' own samples (no weights
involved), the hard read of a fixed reasoning changes under a candidate in only $0.2$--$3.2\%$ of
(candidate, question, sample) triples; the read-off term is $0.1$--$1.2$ accuracy points per edit,
against held-out effects of $3.4$--$10.9$ points, and it does not predict them (Spearman $-0.31$ to
$+0.28$; Table~\ref{tab:readoff}). Two consequences follow. First, a signal computed on a fixed
reasoning measures a small and, on wrongly answered questions, unfaithful part of an edit's effect
(\S\ref{sec:diagnosis} finds exactly this signature). Second, edits move the reasoning
distribution far: the divergence $\mathrm{KL}(\pi_P\,\|\,\pi_{P'})=-\mathbb{E}_{\pi_P}[\log w]$ has a
median of $4.5$--$14.6$ nats per reasoning across pools, so the reasoning a fixed-trace gradient is
taken over is rarely a likely reasoning under the edited prompt.

\paragraph{The natural repair: reweight the model's own samples.} If the missing term is $D$,
the obvious fix keeps \greater{}'s white-box access and estimates $D$ from samples drawn once
under the incumbent. Sample $r_1,\dots,r_S\sim\pi_P^\tau(\cdot\mid x)$ per training question,
read them under the incumbent, and compute for every candidate the exact log-ratios $\log w_s$ of
the \emph{same} samples (teacher forcing, or one superposed pass, \S\ref{sec:calculus}). The
tempered self-normalized importance estimate
\begin{equation}
\widehat{J}_\beta(P'\mid x)=\frac{\sum_s w_s^{\beta}\,c_P(x,r_s)}{\sum_s w_s^{\beta}},\qquad
\widehat{\Delta}_\beta(P')=\mathbb{E}_x\big[\widehat{J}_\beta(P'\mid x)-\tfrac1S\textstyle\sum_s c_P(x,r_s)\big]
\label{eq:snis}
\end{equation}
estimates $D$ (with candidate reads $c_{P'}$ instead of $c_P$ it estimates the full change).
$\beta=0$ recovers the fixed-reasoning family (no reweighting); $\beta=1$ is consistent as
$S\to\infty$; intermediate $\beta$ trades bias for variance, and for small $\beta$,
$\widehat\Delta_\beta\approx\beta\,\widehat{\operatorname{Cov}}(\log w,c)$. Because importance
weights over hundreds of tokens are heavy-tailed, we choose $\beta$ \emph{without labels} as the
largest value in $\{1,\tfrac12,\tfrac14\}$ whose median effective sample size
$(\sum_s w_s^\beta)^2/\sum_s w_s^{2\beta}$ is at least $S/4$. No reasoning is generated and no
answer is read under any candidate. The catch is coverage: importance sampling needs a sample of
size about $\exp(\mathrm{KL})$ between target and proposal \citep{chatterjee2018sample}, and the
divergences measured above are around ten nats per reasoning. Self-normalization hides this---the
weights are renormalized among samples that are all unlikely under the candidate---so the estimate
can look well-conditioned while measuring nothing. We nevertheless test it, preregistered, because
it is the principled ``gradient over the distribution of reasoning'' that \greater{}'s objective
lacks (\S\ref{sec:exp-validity}).

\paragraph{Why verification is noisy.} Greedy correctness is a $\{0,1\}$ quantization of a
per-question success probability. Comparing two prompts on $n$ questions by greedy accuracy has
variance $\approx f/n$, where $f$ is the fraction of questions whose correctness flips between
them; single edits flip $16$--$22\%$ of questions while moving accuracy by about two points
(\S\ref{sec:diagnosis}), so on an $8$-question minibatch the comparison has a standard deviation of
about $15$ points. Selecting the best of many candidates on such a comparison is a winner's curse.
